\chapter{Redes de insumo-producto}

\section{Construcción de la red de insumo-producto internacional}

Los datos de las matrices insumo-producto multirregionales han demostrado ser una herramienta confiable para analizar la globalización económica. Gracias a ellos, se puede estudiar las cadenas globales de valor y estudiar de manera simultánea las interrelaciones en las economías nacionales y regionales \parencite[p.~18]{Xing2022}. Una forma de realizar este estudio es representar la información del modelo de insumo producto como una red y utilizar los métodos de la teoría de redes. 

A partir de la submatriz de consumo intermedio de una \acrfull{mriot} es posible establecer una red del flujo comercial interindustrial internacional. Cada sector de cada país  es representado por un vértice de un grafo, y se establece una arista entre dos vértices si existe flujo comercial entre sus correspondientes sectores. Debido a que el flujo comercial tiene una dirección, del sector vendedor al sector comprador, la arista puede considerarse como dirigida. El valor del flujo comercial de un sector a otro se asocia a su respectiva arista. De esta manera, la interdependencia del sistema económico internacional es modelada como un grafo dirigido y ponderado.

Matemáticamente, un grafo $G$ se define formalmente como un par ordenado $(V,E)$ donde $V$ es un conjunto de vértices y $E$ un conjunto de aristas. El número de vértices $n=|V|$ es denominado el orden del grafo y el número de aristas $m=|E|$ su tamaño \parencite[p.~2]{Bondy2008}. Dos vértices son adyacentes si están conectados por una arista; y una arista que tiene como extremo a un vértice se dice que es incidente a éste. El número de aristas incidentes a un vértice se denomina el grado del vértice. Un grafo puede ser representado por su matriz de adyacencia $A(G) = [a_{ij}]$, donde $a_{ij}=1$ si existe una arista $(i,j)\in E$ y $0$ en caso contrario. Las aristas dirigidas son conocidas como arcos, y al grafo que las contiene como digráfica o grafo dirigido. Un arco es un par ordenado de vértices $e=(i,j), \ i,j\in V$ donde el primer vértice $i$ indica el punto de inicio, y el segundo vértice $j$ es el punto de llegada \parencite[p.~31] {Bondy2008}. Cuando el grafo representa un sistema se denomina red \parencite[p.~7]{Estrada2012}. Un grafo o una red es ponderada, si se asocia un número real o peso $w(e)$ a cada arista $e\in E$, y es denotada con $G=(V,E,w)$ \parencite[p.~50]{Bondy2008}. 

Para construir el conjunto de redes insumo-producto internacional basadas en las tablas \gls{icio} se siguió el siguiente procedimiento:
\begin{itemize}
    \item Para cada tabla anual \gls{icio}, se consideró la submatriz cuadrada de demanda intermedia $\textbf{Z}(t)=\left[z_{ij}^t \right]$, donde $z_{ij}^t$ es el valor de los bienes y servicios que un sector-país $i$ vende al sector-país $j$ en el año $t$, a precios corrientes y del productor, en dólares estadounidenses.\footnote{Debido a su alto nivel de agregación, no se consideraron los datos agrupados como ``resto del mundo'' (\gls{row}).} Todas las submatrices tienen $n=80*50$.
    \item Para cada año se calculó el flujo  intermedio total $tot(t)$, entendido como el valor total de los intercambios comerciales entre todos los sectores de todos los países considerados en el año $t$. Esta cantidad es igual a la suma de los totales por filas o por columnas\footnote{Observe que el sistema no es cerrado, ya se ha descartado la información del valor agregado y de la demanda final. Por lo tanto, la suma de una fila $i$  no necesariamente coincide con la suma de la columna $i$.}:
    \begin{equation}
    	tot(t) = \sum_{j=1}^{n} \left(\sum_{i=1}^{n} z_{ij}^t\right)= \sum_{i=1}^{n} \left(\sum_{j=1}^{n} z_{ij}^t\right) 
    \end{equation}
    \item Después, se calcula la matriz de proporciones interindustriales
	 \begin{equation}
	 	\textbf{Z}^{'}(t) = \left[\frac{z^t_{ij}}{tot(t)}\right] = \left[ x^t_{ij}\right]
	 \end{equation}
    dividiendo cada elemento de $\textbf{Z}(t)$ entre $tot(t)$. Los elementos de esta matriz
    son la proporción de cada flujo interindustrial respecto al total del comercio intermedio de un año específico. De esta forma, se pueden hacer comparaciones históricas sin verse afectadas por las diferencias anuales en los precios corrientes.
    \item Para cada año se construyó una red \gls{icio} $G^{t}=(V^{t}, E^{t}, w^{t})$, tal que:
    \begin{itemize}
        \item El conjunto de vértices $V^{t}$ está constituido por todas las combinaciones de países y sectores económicos, con un total de $n=4,000$ vértices en cada grafo.
        \item Se establece un arco en $E^{t}$ entre dos vértices $i,j\in V^t$ (incluyendo $i=j$) si su correspondiente proporción en $\textbf{Z}^{'}(t)$ es positiva, es decir, si $x^t_{ij}>0$. %No se consideraron los lazos, es decir, las aristas $e=(i,i), \ i\in V$.
        \item La función de pesos $w^t : E^t \rightarrow (0,1)$
        se define como
     	\begin{equation}
        	w^t(i,j) := x^t_{ij}
     	\end{equation}
    \end{itemize} 
\end{itemize}

%En la figura \ref{Fig_RedIcio} se ilustra el proceso de construcción para un modelo con 3 países y 3 sectores.

%\begin{figure}[H]
%\centering
%\begin{subfigure}{.32\textwidth}
  %  \includegraphics[width=\textwidth]{A.png}
%    \caption{\footnotesize{MRIOT\\ $ $}}
%\end{subfigure}
%\begin{subfigure}{.32\textwidth}
  %  \includegraphics[width=\textwidth]{B.png}
%    \caption{\footnotesize{Matriz de proporciones\\interindustriales}}
%\end{subfigure}
%\begin{subfigure}{.32\textwidth}
%    \includegraphics[width=\textwidth]{C.png}
%    \caption{\footnotesize{Red de insumo-producto\\ internacional}}
%\end{subfigure}
%\label{Fig_RedIcio}
%\end{figure}

%Los otros elementos del modelo insumo producto, es decir, los datos de la demanda final y el valor agregado también pueden integrarse en un enfoque basado en redes. En este modelo está formado por tres capas. La primera es formada por vértices que corresponden a los elementos del valor agregado de cada país-sector, y con aristas que apuntan a los vértices de la matriz de uso intermedio.  % Noción de multicapas Xing (2022) p. 19. 

\section{Detección de comunidades}

La detección de comunidades es un problema de la teoría de redes que consiste en buscar una división de los vértices de una red en grupos o comunidades, de tal forma que los vértices dentro de cada comunidad estén densamente conectados entre sí, mientras que los vértices que pertenecen a distintas comunidades estén escasamente conectados \parencite[p.~2]{Blondel2008}.  Las definiciones tradicionales de comunidad comparan la proporción de aristas dentro y entre las comunidades. Sin embargo, el punto de vista moderno se enfoca en la probabilidad de que un vértice comparta aristas con un subgrafo. La existencia de comunidades implica que los vértices interactúan con mayor intensidad con otros miembros de su comunidad, que con los vértices de otras comunidades \parencite[p.~7]{Fortunato2016}.

La función más utilizada para medir la calidad de una partición de los vértices en comunidades es la modularidad \parencite[p.~29]{Fortunato2016},  definida por   \textcite{Newman2004}. Dada una $k$-partición\footnote{Una $k$-partición de un conjunto $V$ se define como una colección de conjuntos $$P=\{C_1, C_2, \ldots, C_k\}$$ que satisface tres condiciones:
	\begin{enumerate}
		\item La unión de todas las partes es igual al todo $$\bigcup_{C_i\in P}C_i=V$$
		\item Las partes son disjuntas por pares $$C_i \cap C_j =\emptyset, \quad \forall C_i,C_j\in P, \ i\neq j$$
		\item No hay partes vacías $$C\neq \emptyset, \quad \forall C_i\in P$$
\end{enumerate}} $P=\{C_1,C_2,\ldots,C_k\}$ de $V$, la modularidad se define como
\begin{equation}
    Q(P) = \frac{1}{2m}\sum_{i,j\in V} \left(a_{ij}-p_{ij}\right) \delta(C_i,C_j) \label{C3_eq_modularity_general}
\end{equation}
donde $a_{ij}$ son los términos de la matriz de adyacencia de $G$; $p_{ij}$ representa el valor esperado de conexión entre los vértices $i$ y $j$, si las conexiones de la red se reconectaran de forma aleatoria, pero respetando ciertas propiedades de la red original; y $\delta(C_i, C_j)$ es una función que es igual a 1 si los vértices $i$ y $j$ pertenecen a la misma partición y 0 en caso contrario, conocida como delta de Kronecker. La elección estándar es suponer que la red aleatoria de comparación (\textit{modelo nulo}) preserva en promedio el grado de todos los vértices. Bajo este supuesto,  la probabilidad de que los vértices $i$ y $j$ con grados $k_i$ y $k_j$ es igual a $\frac{k_i}{2m}$ y $\frac{k_j}{2m}$ respectivamente, por lo tanto, el número de conexiones esperados entre $i$ y $j$ es igual a \parencite[p.~29]{Fortunato2016}:
\begin{equation}
	p_{ij} = 2mp_ip_j = \frac{k_ik_j}{2m}
\end{equation}

Sustituyendo el término $p_{ij}$ en \eqref{C3_eq_modularity_general} se llega a la formulación clásica de la modularidad:
\begin{equation}
	Q(P) = \frac{1}{2m}\sum_{i,j\in V} \left(a_{ij}-\frac{k_ik_j}{2m}\right) \delta(C_i,C_j) \label{C3_eq_modularity1}
\end{equation}

%El término $\frac{k_i}{2m}\frac{k_j}{2m}=p_ip_j$ puede interpretarse como el número de aristas esperados entre $i$ y $j$, si las aristas de la red fueran reconectadas de manera aleatoria, pero conservando en promedio el grado de todos los vértices . 

Debido a que la delta de Kronecker sólo contribuye para los pares de vértices que pertenecen al mismo grupo, la suma puede agruparse por clúster y reformularse\footnote{Los pasos detallados de esta reformulación pueden encontrarse en la sección 9.11 de \textcite{Barabasi}.} de la siguiente manera
\begin{equation}
Q(P) = \sum_{C\in P} \left[\frac{l_C}{m}-\left( \frac{k_C}{2m}\right)^2 \right] \label{C3_eq_modularity2}
\end{equation}
donde $l_C$ es el número total de aristas que unen vértices de la comunidad $C$ y que $k_C$ es la suma de los grados de los vértices que pertenecen a la comunidad $C$. El primer término de cada sumando de la ecuación \ref{C3_eq_modularity2} es la fracción de aristas del grafo que están dentro de la comunidad $C$, mientras el segundo término es la fracción de aristas esperadas en la comunidad $C$ si el grafo se conectara de manera aleatoria, preservando los grados de los vértices en promedio \parencite[p.~29]{Fortunato2016}. 

El valor de la modularidad $Q(P)$ se encuentra acotada dentro del rango $[-1,1]$. Si la proporción de aristas dentro de las comunidades es igual o peor que las esperadas en el modelo aleatorio se obtienen valores de  $Q(P)$ cercanos a 0 o negativos; mientras que una modularidad igual a 1, indicaría una estructura comunitaria muy fuerte; aunque en la práctica, la modularidad en redes con comunidades reales suele toma valores entre  0.3  y 0.7 \parencite[p.~7]{Newman2004}. La figura \ref{Fig_ModularidadEjemplo} muestra un comparación entre la modularidad de una partición que refleja la estructura comunitaria de la red y la modularidad de una partición aleatoria.


\begin{figure}[H]
\centering
\begin{subfigure}{.45\textwidth}
	\centering
  \includegraphics[width=\textwidth]{Modularidad1.png}
\end{subfigure}
\hfill
\begin{subfigure}{.45\textwidth}
	\centering
  \includegraphics[width=\textwidth]{Modularidad2.png}
\end{subfigure}
    \caption{Modularidad de dos particiones de una misma red.}
    \label{Fig_ModularidadEjemplo}
\end{figure}

La modularidad de una comunidad específica $C\in P$ puede evaluarse mediante el término
\begin{equation}
	Q(C) = \frac{l_C}{m}-\left( \frac{k_C}{2m}\right)^2
\end{equation}
el cuál mide la proporción de aristas dentro de una comunidad con la proporción de aristas esperadas en los vértices de la comunidad bajo el modelo nulo. De esta forma, la modularidad de la partición es simplemente la suma de las modularidades de cada clúster:
\begin{equation}
	Q(P) = \sum_{C\in P} Q(C)
\end{equation}


%Conforme mayor sea el valor del primer sumando respecto al segundo, más certeza se tiene que la densidad de aristas de $C$ provienen de una estructura comunitaria y no de una disposición aleatoria. Por lo tanto, se considera que valores grandes de $Q$ indican una estructura comunitaria en la red . \parencite[p.~29]{Fortunato2016}.

%De la ecuación \ref{C3_eq_modularity2} se desprende fácilmente el rango en el que se encuentra la modularidad. Debido a que $\left( \frac{k_C}{2m}\right)^2 \geq 0$ y $0<\sum_C l_C \leq m$ se cumple que
%\begin{equation}Q(P) \leq \sum_C \frac{l_C}{m} =\frac{1}{m} \sum_C l_C \leq \frac{m}{m}=1\end{equation}

La función de modularidad de Newman y Girvan se basa en la la idea de que, en una red conectada de forma aleatoria, sus vértices no se agrupan en clústeres, por lo tanto, la existencia de un agrupamiento se revela mediante la comparación de la densidad real de las aristas en una comunidad y la densidad que se esperaría si los vértices de la comunidad se unieran sin considerar la estructura de comunidad, es decir, de forma aleatoria \parencite[p.~89]{Fortunato2010}

Si la red tiene pesos por aristas, es posible extender la definición de modularidad, remplazando los grados $k_i$ y $k_j$ en la ecuación \ref{C3_eq_modularity1} por los pesos acumulados $s_i$ y $s_j$ de los vértices $i$ y $j$. La fuerza o el peso acumulado de un vértice $i$ se define como la suma de los pesos de las aristas incidentes  al vértice:
\begin{equation}
	s_i = \sum_{j \in N(i)} w(i,j), \qquad v\in V
\end{equation}
donde $N(i)$ representa el conjunto de vecinos o vértices adyacentes de $i$. Para una normalización adecuada el número total de aristas $m$ se sustituye por $W$, la suma de los pesos de todas las aristas. La modularidad para el caso ponderado se define como 
\begin{equation}
	Q_w(P) = \frac{1}{2W} \sum_{i,j\in V}\left(w_{ij} -\frac{s_is_j}{2W} \right)\delta(C_i, C_j)
\end{equation}
donde $w_{ij}$ es el peso de la arista $(i,j)$. Si además la red es ponderada y dirigida, la modularidad se puede calcular como:
\begin{equation}
	Q_{w,d}(P) = \frac{1}{W} \sum_{i,j\in V}\left(w_{ij} -\frac{s_i^{out}s_j^{in}}{W} \right)\delta(C_i, C_j)
	\label{C3_eq_modularity_directed_weighted}
\end{equation}
donde $s_i^{out}$ denota el peso de las aristas que salen de $i$, mientras que $s_j^{in}$ representa el peso de las aristas que llegan a $j$ \parencite[p.~107]{Fortunato2010}.


\subsection{Algoritmo de Leiden y Louvain} \label{LeidenLouvain}

Se ha demostrado que el problema de encontrar una partición de una red que maximice la modularidad pertenece a una clase de problemas conocidos como \textbf{NP-completo} \parencite{Brandes2006}.\footnote{Un problema de decisión, es decir, cuya respuesta es un valor booleano $\{si, no\}$, se considera de la clase \textbf{P} si puede resolverse en tiempo polinomial, esto es, que su tiempo de resolución esté acotado por una función polinómica respecto al tamaño de la entrada. Por otra parte, los problemas de clase \textbf{NP} son aquellos para los cuales, si la respuesta es un \textit{si}, existe una prueba o solución que puede verificarse en tiempo polinomial. Se sabe que cualquier problema en \textbf{P} está también en \textbf{NP}, pero la contención inversa es uno de los siete problemas matemáticos del milenio, aún sin resolver. Se dice que un problema $\Pi$ es \textbf{NP-duro}  si la existencia de un algoritmo que resuelva $\Pi$ en tiempo polinomial implicaría que todos los problemas en \textbf{NP} pueden resolverse también en tiempo polinomial. De forma intuitiva, es un problema al menos tan difícil como cualquier problema en \textbf{NP}. Finalmente, los problemas \textbf{NP-completos} son los problemas \textbf{NP-duros} que a su vez están en la clase \textbf{NP} \parencite[p.~381-383]{Erickson2019}. 
} Por tal razón, se han desarrollado diversos algoritmos heurísticos para encontrar soluciones aproximadas a este problema. En \textcite{Fortunato2010} se encuentra una amplia revisión de estos algoritmos.

En esta tesina, utilizo dos de los algoritmos más prometedores en este campo: el algoritmo de Louvain \parencite{Blondel2008}, y una versión mejorada, el algoritmo de Leiden \parencite{Traag2019}, que toman su nombre por la ciudad en donde trabajan sus autores. Ambos algoritmos fueron diseñados para maximizar la modularidad, incluyendo su definición ponderada y dirigida.
% En la práctica, el algoritmo de Leiden ha mostrado buenos resultados en la calidad de las particiones obtenidas y la eficiencia computacional en pruebas con redes reales (véase \parencite{Traag2019, Aref2024}). 
%\footnote{Aunque el algoritmo original fue definido para maximizar la modularidad, pueden utilizarse otras funciones de calidad de la partición \parencite[p.~1]{Traag2019}.}

Los dos algoritmos se basan en un enfoque multinivel. En la primera fase, comienzan con una partición inicial, la cual van mejorando iterativamente, hasta encontrar una solución que no se puede mejorar (un máximo local). La segunda fase consiste en \textit{agrupar} o \textit{reducir} el grafo, esto es, construir un nuevo grafo cuyos vértices son los clústeres de la partición final de la primera fase; mientras que las aristas y pesos del nuevo grafo, representan la información agrupada de las aristas y pesos del grafo anterior. Estos pasos se repiten alternadamente.

El algoritmo de Louvain se compone de dos fases principales: movimientos locales de vértices y agregación de la red. Usualmente, el algoritmo comienza con una partición inicial en la que cada vértice es una comunidad, aunque es posible comenzar con otro tipo de particiones. Luego, recorre cada vértice del grafo y evalúa si mover el vértice a otra comunidad aumenta la modularidad de la nueva partición; esto incluye considerar el vértice como una nueva comunidad. Si la ganancia de modularidad es positiva, el vértice se traslada a la comunidad que ofrece el mayor aumento. Este proceso se repite hasta que ningún vértice pueda cambiar de comunidad y mejorar la modularidad. A partir de la última partición encontrada en la fase anterior, se construye una red agregada, en la que cada vértice representa una comunidad detectada. En esta nueva red, se establece una arista entre dos comunidades diferentes si, en la red anterior, existía al menos una arista entre un par de vértices pertenecientes a esas comunidades. Si la red es ponderada, el peso de cada arista en la red agregada es igual a la suma de los pesos de todas las aristas que conectaban vértices entre las dos comunidades originales. Estas dos fases se repiten alternadamente hasta que cada comunidad de la partición se conforme de un sólo vértice \parencite{Fortunato2010, Traag2019}.  

En términos matemáticos, la agregación de la red en la $i$ se establece de la siguiente forma: Sea $G^{i}=(V^{i}, E^{i}, w^{i})$ una red ponderada por aristas en la iteración $i$ y $P^{i}=\{C_1, \ldots, C_k\}$ una partición de los vértices de $G^{i}$. Entonces,  la red agregada es $G^{i+1}=(V^{i+1}, E^{i+1}, w^{i+1})$ con
\begin{eqnarray}
V^{i+1} &=& P^i\\
E^{i+1} &=& \{(C, D) \ | \ C, D\in P^i, \ C\neq D, \ (u,v) \in E^{i}, \ u\in C, v\in D\}\\
w(C,D) &=& \sum_{\substack{(u,v)\in E\\ u\in C,\, v\in D}} w(u,v), \quad \forall C,D\in P^{i}, \ C\neq D
\end{eqnarray} 

Se ha demostrado teóricamente que el algoritmo de Louvain puede encontrar comunidades mal conectadas. En particular, dadas ciertas redes, puede conducir a identificar comunidades internamente disconexas \parencite[p.~2]{Traag2019}. El algoritmo de Leiden puede resolver estos problemas. Este algoritmo tiene tres fases: 1) movimiento local de vértices; 2) refinamiento de la partición; y 3) agregación de la red basada en la partición refinada y estableciendo la partición no refinada de la fase 2 como una partición inicial de la red agregada. El algoritmo de Leiden introduce una mejora en la fase de movimiento de vértices que lo hace más rápido. Para ello, se establece una cola de vértices a recorrer, ordenados de forma aleatoria, la cual asegura que todos los vértices se visiten al menos una vez, pero reduce el número de veces que un mismo vértice es seleccionado. La idea de la fase de refinamiento consiste en que dada una partición $P$ obtenida en la fase 1, sus comunidades son divididas en múltiples subcomunidades. La partición refinada obtenida en la fase 2 sirve para crear la red agregada, mientras que la partición no refinada sirve como partición inicial de esta nueva red. Las tres fases se suceden alternadamente. El algoritmo de Leiden es considerablemente más complejo que el algoritmo de Louvain. Sin embargo, el algoritmo de Leiden tiene propiedades teóricas que aseguran la buena conexión de las comunidades; además, de obtener mejores resultados en términos de calidad y tiempo de ejecución en redes reales (véase \textcite{Traag2019, Aref2024}).

Debido a que estos dos algoritmos tienen un componente aleatorio al realizar la busqueda de movimientos, no es seguro que lleguen a la misma partición en diferentes ejecuciones, incluso para la misma solución inicial. Por eso es recomendable ejecutar los algoritmos varias veces y evaluar la calidad y variabilidad de sus soluciones. % Por ello es recomendable ejecutar cada algoritmo varias veces fijando la semilla que genera números pseudoaleatorios en la impletación computacional. 

%\subsection{Implementación computacional}

Para la implementación computacional de los dos algoritmos se utilizó la librería \texttt{leidenalg} de \texttt{Python} \parencite{Traag2023}, la cual fue desarrollado por Vicent Traag, uno de los autores del algoritmo de Leiden. La librería fue programada en el lenguaje \texttt{C++}, lo cual asegura una gran rapidez de ejecución, pero pueden importarse las funciones y algoritmos en \texttt{Python}. Esta librería permite ejecutar tanto el algoritmo de Louvain como el de Leiden para los casos dirigidos y ponderados.

\section{Detección de comunidades en redes de flujos comerciales}

La detección de comunidades ha sido aplicada desde hace dos décadas como una herramienta para analizar la estructura de las redes de flujos comerciales internacionales. Principalmente, los estudios se han enfocado en identificar la formación de comunidades económicas, sin depender de factores exógenos, como la cercanía geográfica o los vínculos institucionales de integración económica, sino únicamente de la densidad de relaciones económicas (véase \textcite[p.~191]{Tsekeris2017}).

Inicialmente, estos métodos fueron utilizados para analizar la red de comercio mundial (\textit{International Trade Network}, \textit{World Trade Network}), en donde los vértices corresponden a países, y las aristas y sus pesos indican los flujos de importación y exportación entre ellos \textcite{Tzekina2008, Piccardi2012, Reyes2014, Fan2014, Zhu2014, Antonietti2022, Bartesaghi2022}. Sin embargo, el tipo de grafo y la definición de los pesos varía de un trabajo a otro. En \textcite{Tzekina2008, Piccardi2012} consideran grafos dirigidos, mientras que en \textcite{Zhong2014, Fan2014, Reyes2014, Antonietti2022, Bartesaghi2022} las aristas no son dirigidas, aunque en cada artículo tienen una forma particular de ponderar las aristas.  Por ejemplo, en \textcite{Piccardi2012} los pesos de una arista son simplemente el valor de las exportaciones de un país a otro; en \textcite{Zhong2014} los pesos equivalen a la suma de importaciones y exportaciones, es decir, el flujo comercial bilateral total; mientras que en \textcite{Tzekina2008} dividen el flujo bilateral de cada arista entre el volumen de comercio total del país exportador. Respecto a los métodos de detección de comunidades, solo en \textcite{Piccardi2012, Reyes2014, Zhu2014, Fan2014} buscan optimizar la modularidad de las particiones mediante algoritmos heurísticos.
%mientras que en \textcite{Reyes2014} el grafo es no dirigido y las aristas son ponderadas por la media geométrica del flujo comercial entre dos países.
% ITN-WTN
% Tzekina 2008 - ponderada dirigida WTW - umbral - 85 
% Reyes 2014 - ponderada no dirigida - Modularidad espectral -62
% FAN 2014 - ponderada no dirigida - Nodularidad WEO - 3 comunidades principales
% Antonietti - Ponderada no dirigida - communicability distance umbral
% Bartesagyu, 2022 - ponderada no dirigida - communicability umbral
% Especific commodity
% Barigozzi 2010 - ponderada y dirigida WTW - modularidad Tabu 

El trabajo de \textcite{Barigozzi2011} amplia el uso de detección de comunidades a la red de comercio internacional desagregada por productos. Consideran 192 países y se enfocan en los 14 productos más comercializados. Sus grafos son dirigidos y el peso de sus aristas es igual al valor monetario de las exportación de un producto específico entre dos países. Para cada año y producto definen particiones utilizando el algoritmo de búsqueda tabú para optimizar la modularidad. Comparan éstas particiones con las que se identifican en la red agregada, así como con comunidades preestablecidas definidas a partir de acuerdos comerciales. %Encuentran que las particiones de la red agregada, generalmente, tiene menor número de comunidades que las particiones de las redes de un producto específico, lo cual sugiere 
Otros estudios han analizado la estructura comunitaria de redes de comercio internacional de productos específicos, por ejemplo, la red de comercio de petróleo \parencite{Zhong2014}, de cobre \parencite{Dong2018}, de alimentos \parencite{Torreggiani2018}, de metales resiclables \parencite{Hu2020} y de pesticidas \parencite{Li2023}. Recientemente, en  \textcite{Tsekeris2025} analizaron la evolución de las comunidades en la red de comercio de gas, acotada a países de Europa, el norte de África y Medio Oriente, en el contexto de la guerra entre Rusia y Ucrania.

Las redes de comercio internacional se basan en datos de importaciones y exportaciones entre países, agrupando tanto el comercio intermedio, como aquel que es destinado a satisfacer la demanda final. Por lo tanto, su estructura comunitaria puede reflejar el grado de agrupamiento comercial que hay entre los países, pero no ayuda a develar los vínculos de interdependencia que se forman a nivel productivo. Algo que caracteriza la estructura comunitaria identificada en la mayoría los trabajos mencionados es que el número de comunidades es muy limitado, generalmente en un rango de entre 2 y 6 comunidades, pese a que el número de países considerados supera los 160 (véase \textcite[p.~2058]{Barigozzi2011}, \textcite[p.~3]{Piccardi2012}, \textcite[p.~73]{Fan2014} y \textcite[p.~3]{Zhu2014}). %En los casos donde el número de comunidades es mayor, ocurre que sólo hay dos o tres comunidades con varios países, y las demás comunidades están constituidas por un sólo país (véase \parencite{Bartesaghi2022}). 
Esto ha llevado a algunos autores a sugerir que la red de comercio internacional carece de una estructura comunitaria claramente definida, lo cual constituye evidencia de que la integración económica se manifiesta como un proceso de globalización más que de regionalización \parencite[p.~10--11]{Piccardi2012}
% Barigozzi2011 - 162 países y 2-4 comunidades
% Piccardi 125, 181 paísesy 2-4 comunidades
% Zhou  200 - 3 comunidades

% Autor que menciona la agregación como algo importante.

A diferencia de la red de comercio mundial, la red de insumo-producto desagrega los intercambios comerciales según las industrias involucradas, lo que permite analizar simultáneamente los vínculos de interdependencia productiva tanto a nivel nacional como internacional. Esta representación ha sido empleada en diversos estudios que aplican métodos de detección de comunidades para identificar la formación y evolución de bloques económicos. Estos bloques se definen como agrupaciones de sectores altamente cohesionadas e integrados en función de su elevada densidad relativa de flujos interindustriales. Los trabajos que han adoptado  este enfoque son:
\begin{itemize}
	\item En el trabajo pionero de \textcite{Cerina2015} se analizan las características de la red de insumo-producto mundial a nivel global, regional y local. A nivel global, los autores calculan varias métricas de la topología de la red y concluyen que las industrias están conectadas de forma asimétrica, lo cual implica que pequeños choques económicos pueden provocar fluctuaciones macroeconómicas significativas. Para el análisis regional, aplican un enfoque de detección de comunidades y encuentran que la mayoría de las comunidades identificadas corresponden a países o grupos de países cuyos sectores están completamente integrados dentro de una misma comunidad. Los autores argumentan que esto es evidencia de que la producción mundial continúa operando a escala nacional o regional. En particular, señalan el surgimiento de un gran bloque económico europeo liderado por Alemania; mientras que para Asia no se obtienen conclusiones debido al limitado número de economías asiáticas incluidas en los datos. A nivel local, comparan diferentes medidas de centralidad de los vértices de la red con el fin de identificar sectores clave que cumplen un rol fundamental en la red. 
	% El análisis se basa en redes anuales construidas a partir de los datos de WIOD para el periodo 1995-2011, que cubre 40 países y 35 industrias. Todas las redes son dirigidas y ponderadas, y los pesos de las aristas representan flujos monetarios interindustriales a precios corrientes del productor (en millones de dólares estadounidenses), sin normalización. No se incluyen los datos del valor agregado ni de la demanda final. La detección de comunidades se realiza mediante un algoritmo heurístico jerárquico para maximizar la modularidad, propuesto en \parencite{Newman2004}.
	
	\item En \textcite{Maluck2015} comparan las particiones encontradas con métodos de detección de comunidades en la red de insumo-producto global\footnote{Aunque, sus datos corresponde a una \gls{mriot} que incluye relaciones interindustriales, los autores utilizan en término \textit{red de comercio internacional}.} con particiones preestablecidas por país o sector, 
	%utilizan la base de datos Eora, con 186 países y 26 sectores, para construir redes insumo-producto mundial para los años 1990 a 2011. Consideran el caso dirigido y no dirigido, asignando como peso de arista el valor monetario del flujo comercial normalizado por el volumen de comercio global anual. Utilizan el algoritmo de Louvain para encontrar particiones en cada red anual.
	es decir, que agrupan todos los vértices correspondientes aun país/sector en una misma comunidad. La modularidad de las particiones detectadas indica una estructura comunitaria bien definida (entre 0.68 y 0.76), mientras que la modularidad de la partición basada en el país es similar y ligeramente menor durante todo el periodo. Por el contrario, la modularidad de la partición basada en sectores siempre se encuentra por debajo de 0.05, mostrando que la integración de un mismo sector entre países, no refleja la estructura de integración de los flujos interindustriales.

	
	\item En \textcite{Tsekeris2017}, el objetivo es analizar los cambios estructurales en las cadenas de valor global,
	% durante el periodo 1995-2011,
	incluyendo la evolución de sus comunidades. 
	%El autor define la red de la cadena de valor país-sector como una red en la que cada vértice representa un par sector-país y las aristas corresponden a los vínculos interindustriales, ponderadas por el valor monetario de las transacciones. Sus datos provienen de las tablas ICIO de la OECD, que incluyen 34 sectores y 62 países más el resto del mundo.
	%Los resultados de la detección de comunidades coinciden con los hallazgos reportados en \parencite{Cerina2015}, en tanto los sectores de un país casi siempre pertenecen a la misma comunidad.
    Sus resultados muestran un decrecimiento en el número de comunidades, que pasa de 26 comunidades en 1995 a 18 en 2011, lo que sugiere que las cadenas de valor nacionales han incrementado su apertura y aumentado su integración regional durante ese periodo \parencite[p.~196]{Tsekeris2017}.
    % Por otra parte, se reporta una disminución de la modularidad, de 0.749 en 1995 a 0.710 en 2011. 
    La comunidad más grande encontrada en 2011, en términos del número de sectores que la componen, corresponde al sudeste asiático, seguido de 3 comunidades formadas por países europeos. Además, el autor sugiere que los resultados indican una menor interdependencia entre comunidades, en comparación con una mayor integración de las cadenas de valor nacionales dentro de las comunidades \parencite[p.~197]{Tsekeris2017}. 
    %Para la detección de comunidades, se empleó un algoritmo heurístico basado en métodos espectrales propuesto en \parencite{Clauset2004} para optimizar la modularidad.

	\item En \textcite{Piccardi2018} definen cuatro tipos de redes insumo-producto mundial para cada año, a partir
	%utilizan un método alternativo a la maximización de la modularidad, basado en caminatas aleatorias, para estudiar la estructura de comunidades de la red insumo-producto mundial entre los años 1995 y 2011. Para su análisis construyen redes ponderadas y dirigidas anuales con datos de WIOD, al igual que en \parencite{Cerina2015}; sin embargo, a diferencia de ese trabajo, aquí los autores definen cuatro redes distintas para cada año,
	a partir de las combinaciones que surgen de incluir o no la demanda final por país como vértice, y de considerar las direcciones de acuerdo con el flujo material o el flujo monetario de las transacciones. Los resultados obtenidos concuerdan con los hallazgos anteriores, en tanto que las comunidades se forman en torno a las economías completas de un país o grupo de países. Sin embargo, destacan que, en algunos casos, hay sectores altamente integrados entre países. Por ejemplo, el sector de equipos de transporte en varios países aparece integrado dentro de la comunidad de Europa central encabezada por Alemania; mientras que la industria textil y del cuero parecen integrarse en la comunidad formada por la economía italiana. Concluyen que cuando se consideran los datos de la demanda final las comunidades tienden a reflejar las economías nacionales, mientras que cuando esta información es excluida, se observan mayores niveles de integración regional.

	
	\item En \textcite{Xu2019} analizan la red insumo-producto mundial del año 2009, 
	%utilizan los datos de WIOD del año 2009, y la misma representación de red que en \parencite{Cerina2015},
	para analizar medidas de centralidad de vértices y aristas, identificar comunidades y extraer lo que denominan la \textit{columna vertebral} de la economía mundial, mediante la eliminación de aristas menos importantes.
	% Respecto a la detección de comunidades, utilizaron un algoritmo heurístico voraz (\textit{greedy}), propuesto en \parencite{Newman2004fast}, para encontrar particiones que maximicen la modularidad. 
	A partir de su análisis identificaron 72 comunidades para la red de 2009, aunque 50 de estas comunidades están conformadas por un sólo sector. La comunidad más grande detectada, en número de sectores, corresponde a Europa central, encabezada por la economía alemana.
	

	\item El artículo de \textcite{Kireyev2022} introduce nuevos métodos para evaluar la robustez de la detección de comunidades aplicadas a la red de insumo-producto mundial, además de proponer diferentes enfoques para identificar los sectores clave dentro de cada comunidad. Para la construcción de la red insumo-producto mundial, consideran tanto la submatriz de demanda intermedia como la demanda final. 
	%De esta manera, el conjunto de vértices incluye tanto a los pares sector-país como a la demanda final de cada país. 
	%Los datos utilizados provienen de la WIOD versión 2016 cubriendo el periodo 2000-2014, con 43 países (más el resto del mundo) y 56 sectores. No se menciona ninguna normalización de los flujos comerciales. 
	Para detectar comunidades utilizan el algoritmo de Louvain, el cual se ejecuta varias veces para cada red para evaluar la estabilidad de las particiones. Además, definen tres configuraciones iniciales para inicializar el algoritmo: 1) una partición trivial, donde cada comunidad contiene sólo un vértice; 2) partición por país, donde todos los sectores de un país pertenecen a la misma comunidad y 3) una partición obtenida el año anterior. Encontraron diferencias significativas en la calidad de las particiones encontradas dependiendo de la configuración inicial utilizada.
	%Encontraron que mientras las particiones en la tercera configuración permanecían casi sin cambios de un año a otro, cuando se utilizaba la partición trivial como solución inical la estructura comunitaria variaba significativamente de un año a otro.
	Reportan valores altos de modularidad, en el rango de  0.76 y 0.80, lo cual indica una fuerte estructura comunitaria.
	% Los resultados computacionales muestran una tendencia creciente en los valores máximos de la modularidad en el periodo entre 2000 y 2009, teniendo su valor máximo en ese último año, a partir del cuál comienza una tendencia decreciente.
	 Sus conclusiones refuerzan los hallazgos encontrados anteriormente de que las comunidades se forman por los sectores de un país o grupo de países. En particular, encontraron que la asignación de sectores de países relativamente pequeños es muy inestable, cambiando de una comunidad a otra cada año, como los casos de Rumania, Malta, Chipre y Grecia. Asimismo, los autores señalan un proceso de crecimiento de la comunidad económica centrada en Alemania, con el aumento de economías de Europa central y del este, y un debilitamiento del bloque económico centrado en la economía rusa.
\end{itemize}

En la tabla \ref{tab_DeteccionMRIOT} se muestra la base de datos, la cobertura temporal, geográfica y sectorial de cada estudio, así como el algoritmo de detección de comunidades utilizado en los estudios de estructura comunitaria de la red insumo-producto mundial. En todos los artículos utilizan redes dirigidas y ponderan las aristas con los flujos monetarios (a precios corrientes) entre sectores, exceptuando en \textcite{Maluck2015} donde normalizan estos valores dividiendolos por el volumen de comercio global anual.


\begin{table}[h]
	\centering
	\caption{Características de los estudios de detección de comunidades en redes insumo-producto globales}
	\label{tab_DeteccionMRIOT}
	\resizebox{\columnwidth}{!}{%
	\begin{tabular}{@{}ccccc@{}}
		\toprule
		\textbf{Base de datos} & \textbf{Período} & \textbf{Países} & \textbf{Sectores} & \textbf{Algoritmo utilizado} \\ \midrule
		WIOD     & 1995-2011 & 40  & 35 & Jeraquico aglomerativo \parencite{Newman2004} \\
		Eora     & 1990-2011 & 186 & 26 & De Louvain \parencite{Blondel2008} \\
		ICIO-OECD& 1995-2011 & 62  & 34 & Espectral \parencite{Clauset2004}\\
		WIOD     & 1995-2011 & 40  & 35 & Caminatas aleatorias \parencite{Piccardi2011}\\
		WIOD     & 2009      & 40  & 35 & Voraz \parencite{Newman2004fast} \\
		WIOD     & 2000-2014 & 43  & 56 & De Louvain \parencite{Blondel2008}\\ \bottomrule
	\end{tabular}%
	}
\end{table}

Todos estos trabajos muestran que la red de insumo-producto mundial presenta una estructura fuertemente comunitaria, donde los sectores están altamente integrados en economías nacionales, y en algunos casos, en bloques económicos regionales. Incluso hay coincidencia en algunas bloques regionales detectados, como el caso de la comunidad económica conformada por países de Europa central y oriental, alrededor de la economía de Alemania. Sin embargo, en otros casos no hay consistencia entre estudios respecto a los bloques regionales. Por ejemplo, en el caso de la región de Norteamérica (México, Estados Unidos y Canadá), tanto \textcite{Tsekeris2017} como \textcite{Piccardi2018} identifican una comunidad que integra a todos los sectores de los tres países para los datos de 2011; mientras que \textcite{Kireyev2022} reporta comunidades independientes por país para todo el periodo. Por su parte, \textcite{Cerina2015, Xu2019} encuentran también una comunidad por país, pero muestran que algunos sectores de México y Canadá se integran en el bloque estadounidense, para el año 2011 y 2009 respectivamente. Aunque los estudios utilizan datos provenientes de una misma fuente, como en \textcite{Cerina2015, Piccardi2018, Xu2019}, las diferencias en los bloques económicos identificados responden a la aplicación de distintos métodos de detección de comunidades. Esto muestra la importancia de trabajar con métodos de detección de comunidades que aseguren encontrar particiones de mejor calidad.

% Cerina. 1995 Mexico Usa y canadá separados cada uno, USa tienen y7 mexicanos y 1 canadiense. En 2003 Mexico, Candaá y parte de Usa estan integrados. EN 2011 Parte de Méxic, está integrado aUsa, y dos sectores de cadana. Pero México tiene su propia comunidad.
% Tsekeris 1995 Mexico canadá y EU separados, en 2011 Integrados en una comunidad
% Xu - Canada aparte, Mexico aparte. USA tiene 2 de méxico y 1 de canadá (hay una comunidad con 2 de USa y 1 de canada)


La detección de comunidades también se ha utilizado en la literatura para analizar la estructura de las redes insumo-producto nacionales. Por ejemplo, en \textcite{Mcnerney2013} analizan la estructura comunitaria de las redes de insumo-producto nacionales de 25 países, para los años 2002 y 2006, y caracterizan la distribución del peso por arista de estas redes, así como, la distribución de la fuerza por vértice. 
%Sus datos provienen de las tablas IOnacionales de la OECD de los años 2002 y 2006. 
Este artículo complementa los resultados presentados en el reporte técnico de \textcite{McNerney2009} donde se analizan los datos de sólo 20 países con datos de 200. En \textcite{Tsekeris2017Greek} identifican comunidades sectoriales en la red de insumo-producto de Grecia en 2010 y analizan el patrón de influencia de una comunidad sectorial sobre las demás comparando la densidad relativa de los vínculos entre comunidades. En \textcite{Wang2021} construyen una red insumo-producto regional de China de 2007 y 2012, donde los vértices representan pares sector-unidad provincial, y las aristas sus vínculos interindustriales. Sus resultados muestran que la estructura comunitaria es regional, esto es, que la mayoría de los sectores de una provincia pertenecen a la misma comunidad. Argumentan que esto se debe a que por razones geográficas, históricas y administrativas, los vínculos intra provincia son naturalmente más fuertes que entre provincias. Sin embargo, también observan que las industrias más grandes, como la minería, la extracción de petroleo, gas y minerales, usualmente forma comunidades de un sólo vértice, independiente de la comunidad de la provincia en la que se encuentran dichos sectores. 


A partir de la revisión de literatura, se tomaron algunas medidas metodológicas respecto a la construcción de la red:
\begin{itemize}
	\item En primer lugar, se determinó no considerar la información de los flujos comerciales de un sector de un país consigo mismo. La inclusión de estos flujos internos puede reducir la resolución de los métodos que buscan maximizar la modularidad, debido a que aumenta el valor de $W$ y decrece el término de \textit{penalización} $\frac{s_i^{out}s_j^{in}}{W}$ del modelo nulo. Esto dificulta identificar comunidades pequeñas, ya que favorece la fusión entre comunidades al hacer más fácil que un vínculo entre dos industrias supere el término de \textit{penalización} \parencite[p.~6434]{Mcnerney2013}. 
	\item En segundo lugar, la demanda final no se incluye, siguiendo el criterio de \textcite{Cerina2015}. La demanda final suele tener vínculos más fuertes con la oferta nacional que con la producción internacional; por ello, su inclusión podría indicir la formación de comunidades puramente nacionales y ocultar patrones de integración regional, como advierte los resultados de  \textcite[p.~17--18]{Piccardi2018}.
	\item Finalmente, se decidió normalizar los flujos con respecto al volumen total de comercio intermedio, como en \textcite{Maluck2015}, con el fin de permitir la comparación de cambios temporales en los flujos de la red. Cabe mencionar, sin embargo, que en las pruebas computacionales no se observaron diferencias significativas en la modularidad máxima detectada entre las red con pesos normalizados y aquellas sin normalizar.
\end{itemize}
%	\item En \parencite{Mcnerney2013} estudian la estructura comunitaria de las relaciones inteindustriales en 45 países de la OECD, con un nivel de desagregación de 41 (2002) y 48 (2006) industrias. Aplican varias veces un algoritmo heurístico basado en recocido simulado para encontrar diferentes particiones en cada red, y posteriormente ensamblan esas particiones en una sola. Además, caracterizan la distribución de flujo monetario entre industrias (distribución del peso por aristas), así como la distribución del flujo que pasa por cada vértice (distribución de la fuerza por vértice).	Este artículo complementa los resultados presentados en el reporte técnico de \parencite{McNerney2009} donde analiza la red de insumo-producto de 20 países de la OECD con 39 industrias por país. % flujos de bienes de un sector a otro. 
%En \parencite{Mcnerney2013} construyen redes de consumo intermedio nacionales para 45 países utilizando datos de 2002 y 2006 de la OECD. La función de pesos considerada son los flujos monetarios interindustriales normalizados por la producción bruta más las importaciones y los subsidios.	%\item \textcolor{red}{Wang et al (2021) Regional and sectoral structures of the Chinese economy: A network perspective from multi-regional input-output tables} Datos  de MRIOT de China 2007, 2010 y 2012. Cada vértice es un sector dentro de una provincia. Cada arista su ponderación es el multiplicador de 10 000 CNY en volumen de transacciones. Numero de vértices n=900.


%En los últimos 20 años, se han llevado a cabo varios estudios donde la interdependencia industrial es modelada como una red compleja  \parencite[p.~7]{Xing2022}. Algunos ejemplos de estos trabajos son:
%\begin{itemize}
	%	\item En \parencite{Acemoglu2012} utilizan las tablas de requerimientos técnicos\footnote{Los coeficientes de requisitos directos muestran la cantidad de insumos comprados directamente para producir un dólar de cierto bien o servicio \parencite{BEA2005}.} producto-por-producto de los Estados Unidos para el periodo 1972-2002 para construir una red dirigida y ponderada. En este caso los productos corresponden a los vértices y los valores nominales de la matriz los pesos por aristas.	
	%	\item En \parencite{Contreras2014} construyen la red consumo intermedio  para los países de la Unión Europea utilizando las tablas IO de la Eurostat del año 2005. Sólo consideran los flujos positivos para establecer las conexiones entre sectores. No hay mención de un proceso de normalización, por lo que se asume que utilizan los valores nominales de los flujos interindustriales.

	
% \item Río-Chanona (2017). Una red por país y otra por sector WIOD 1995-2009
%\item \textcolor{red}{Bloch, et al. 2011 Vertex centralities in input-output networks reveal the structure of modern economies.} Usan cuentas insumo producto de la base de datos STAN de la OECD, con 47 sectores y 37 países para la versión de 2000. Aunque las tablas reportan monedas locales, solo consideran "unit-free transition matrices"
% \item Grazzini, Spelta (2022) An empirical analysis of the global input-output network and its evolution. Generan un modelo y lo calibran con datos reales. Datos WIOT 43 países y 56 sectores, 2000-2014. Cada vértice es un sector-país, aristas son el flujo nominal de bienes entre sectores
%\end{itemize}

